\chapter{Normalization Methods}

\section{Introduction}

Normalization layers are a central component of modern deep networks. They improve training stability, accelerate optimization, and enable deeper architectures. However, their role is often described informally (e.g., ``reducing internal covariate shift''), which can obscure their true effect.

This chapter develops normalization methods from both mathematical and practical perspectives. We study batch normalization, layer normalization, and related variants, and interpret normalization as a form of reparameterization that improves conditioning and gradient flow.

\section{Why Normalization Helps}

Deep networks suffer from unstable signal propagation and poor conditioning. Activations can grow or shrink across layers, and gradients may become ill-scaled.

Normalization addresses these issues by:
\begin{itemize}
    \item controlling the scale of activations,
    \item improving gradient conditioning,
    \item making optimization less sensitive to parameter initialization,
    \item introducing a mild regularization effect.
\end{itemize}

Rather than viewing normalization as fixing data distributions, it is more accurate to view it as improving the geometry of the optimization problem.

\section{Batch Normalization}

Batch normalization (BN) normalizes activations across a mini-batch.

\subsection{Transformation}

Given activations $x_1,\dots,x_m$ in a batch, define:
\[
\mu_B = \frac{1}{m}\sum_{i=1}^m x_i,
\qquad
\sigma_B^2 = \frac{1}{m}\sum_{i=1}^m (x_i - \mu_B)^2.
\]

The normalized activations are:
\[
\hat{x}_i = \frac{x_i - \mu_B}{\sqrt{\sigma_B^2 + \varepsilon}}.
\]

A learned affine transformation is then applied:
\[
y_i = \gamma \hat{x}_i + \beta,
\]
where $\gamma$ and $\beta$ are trainable parameters.

\subsection{Interpretation}

BN standardizes activations within each batch, then allows the network to recover any necessary scale and shift through $\gamma$ and $\beta$.

This can be interpreted as:
\begin{itemize}
    \item enforcing a normalized intermediate representation,
    \item while preserving expressive flexibility.
\end{itemize}

\subsection{Train/Test Discrepancy}

During training, BN uses batch statistics $(\mu_B, \sigma_B^2)$.  
At test time, it uses running averages:
\[
\mu_{\mathrm{test}}, \quad \sigma^2_{\mathrm{test}}.
\]

\begin{itemize}
    \item this introduces a discrepancy between training and inference,
    \item small batch sizes can lead to noisy estimates.
\end{itemize}

Thus BN depends on batch size and training dynamics.

\subsection{Internal Covariate Shift (Nuanced View)}

The original motivation for BN was to reduce ``internal covariate shift,'' meaning changes in layer input distributions during training.

However, modern understanding suggests:
\begin{itemize}
    \item the main benefit is improved optimization geometry,
    \item BN smooths the loss landscape and stabilizes gradients,
    \item it enables larger learning rates.
\end{itemize}

Thus its effectiveness is better explained by conditioning rather than distributional stability alone.

\section{Layer Normalization and Sequence Models}

Batch normalization depends on batch statistics, which can be problematic for:
\begin{itemize}
    \item small batch sizes,
    \item sequential or autoregressive models.
\end{itemize}

\subsection{Layer Normalization (LN)}

Layer normalization normalizes across features for a single example:
\[
\mu = \frac{1}{d}\sum_{j=1}^d x_j,
\qquad
\sigma^2 = \frac{1}{d}\sum_{j=1}^d (x_j - \mu)^2.
\]

Then:
\[
\hat{x}_j = \frac{x_j - \mu}{\sqrt{\sigma^2 + \varepsilon}},
\qquad
y_j = \gamma \hat{x}_j + \beta.
\]

\subsection{Properties}

\begin{itemize}
    \item independent of batch size,
    \item consistent between training and testing,
    \item widely used in transformers and sequence models.
\end{itemize}

\subsection{Instance and Group Normalization}

Other variants include:
\begin{itemize}
    \item \textbf{Instance normalization:} normalize each channel independently per example,
    \item \textbf{Group normalization:} normalize over groups of channels.
\end{itemize}

These provide flexibility between BN and LN, especially in vision tasks.

\section{Normalization as Geometry and Conditioning}

Normalization can be understood as a reparameterization of the model.

\subsection{Reparameterization View}

Consider a layer output:
\[
z = W x.
\]

With normalization:
\[
\hat{z} = \frac{z - \mu}{\sigma}.
\]

This effectively rescales directions in parameter space, making optimization less sensitive to parameter magnitudes.

\subsection{Conditioning Intuition}

Poorly scaled features lead to ill-conditioned optimization problems. Normalization:
\begin{itemize}
    \item makes level sets more isotropic,
    \item improves gradient alignment,
    \item allows larger and more stable updates.
\end{itemize}

Thus normalization acts as a preconditioner embedded within the network.

\section{When Normalization Fails or Is Unnecessary}

Normalization is not universally beneficial.

\subsection{Failure Modes}

\begin{itemize}
    \item very small batch sizes degrade BN performance,
    \item mismatch between training and inference statistics,
    \item added computational overhead.
\end{itemize}

\subsection{Alternatives}

Some architectures reduce reliance on normalization by:
\begin{itemize}
    \item careful initialization,
    \item residual connections,
    \item architectural constraints.
\end{itemize}

In such cases, normalization may be less critical.

\section{Worked Comparison: BN vs LN}

\begin{itemize}
    \item BN normalizes across the batch:
    \begin{itemize}
        \item depends on batch statistics,
        \item effective in large-batch vision models.
    \end{itemize}
    \item LN normalizes across features:
    \begin{itemize}
        \item independent of batch size,
        \item preferred in sequence and transformer models.
    \end{itemize}
\end{itemize}

Thus the choice depends on data structure and training setup.

\section{A Unifying Perspective}

Normalization layers control scale and improve conditioning within deep networks. Rather than merely stabilizing distributions, they reshape the optimization landscape through reparameterization.

Different normalization methods correspond to different choices of:
\begin{itemize}
    \item normalization axis (batch, feature, channel),
    \item statistical estimation (batch vs per-sample),
    \item invariance properties.
\end{itemize}

Understanding these choices clarifies when each method is appropriate.

\section{Summary}

In this chapter we studied normalization methods in deep learning. We defined batch normalization and analyzed its transformation, training/inference behavior, and practical effects. We introduced layer normalization and related variants, highlighting their role in sequence models.

We interpreted normalization as a geometric reparameterization that improves conditioning and gradient flow. Finally, we discussed when normalization may fail or be unnecessary. These ideas are essential for designing stable and efficient deep architectures.